Une fois le backlog du produit finalisé, nous avons organisé la réunion de planification de sprint. L'objectif de cette réunion était de construire le backlog de sprint en se basant sur le backlog de produit défini par le Product Owner. À l'issue de cette réunion, nous avons identifié les durées prévisionnelles du travail à effectuer pour chaque sprint. Pour ce projet, nous avons divisé le travail en deux releases. La première release contient un seul sprint d'une durée de 2 semaines, tandis que la deuxième release est composée de deux sprints, chacun d'une durée de 3 semaines. La figure 2.2 illustre la répartition des sprints pour notre système.

Décrivez le plan des sprints :

\begin{itemize}
    \item \textbf{Sprint 1 : Collecte de données basée sur caméra}
    \begin{itemize}
        \item \textbf{Objectif} : Mettre en place un système de collecte de données à partir des machines anciennes en utilisant des caméras et des techniques de reconnaissance optique de caractères (OCR).
        \item \textbf{Deliverables} :
            \begin{itemize}
                \item Intégration de la caméra avec le système.
                \item Détection et extraction des valeurs vitales via OCR.
                \item Validation des données collectées par rapport aux données réelles.
            \end{itemize}
        \item \textbf{Durée} : 2 semaines.
        \item \textbf{Tâches clés} :
            \begin{itemize}
                \item Configuration des caméras et des logiciels OCR (OpenCV, Tesseract).
                \item Développement du pipeline de traitement d'images.
                \item Tests de validation des données collectées.
            \end{itemize}
        \item \textbf{Risques} : 
            \begin{itemize}
                \item Mauvaise qualité d'image affectant la précision de l'OCR.
                \item Solution : Utilisation de caméras haute résolution et prétraitement des images.
            \end{itemize}
    \end{itemize}

    \item \textbf{Sprint 2 : Collecte de données à partir des moniteurs Mindray}
    \begin{itemize}
        \item \textbf{Objectif} : Intégrer les moniteurs Mindray pour collecter des données en temps réel via des API ou des protocoles HL7.
        \item \textbf{Deliverables} :
            \begin{itemize}
                \item Connexion des moniteurs Mindray au système central.
                \item Collecte et stockage des données en temps réel.
                \item Validation de la collecte de données via API/HL7.
            \end{itemize}
        \item \textbf{Durée} : 3 semaines.
        \item \textbf{Tâches clés} :
            \begin{itemize}
                \item Configuration des API ou des protocoles HL7 pour les moniteurs Mindray.
                \item Développement du backend pour la collecte et le stockage des données.
                \item Tests de validation des données collectées.
            \end{itemize}
        \item \textbf{Dépendances} : Les données collectées via caméra (Sprint 1) doivent être validées avant de passer à ce sprint.
        \item \textbf{Risques} :
            \begin{itemize}
                \item Problèmes de compatibilité avec les API ou les protocoles HL7.
                \item Solution : Documentation technique et tests approfondis avant l'intégration.
            \end{itemize}
    \end{itemize}

    \item \textbf{Sprint 3 : Système de gestion des alertes et de notification}
    \begin{itemize}
        \item \textbf{Objectif} : Développer un système de gestion des alertes basé sur des seuils prédéfinis et envoyer des notifications au personnel médical.
        \item \textbf{Deliverables} :
            \begin{itemize}
                \item Définition des seuils d'alerte pour les valeurs vitales.
                \item Mise en place d'un système de priorisation des alertes.
                \item Intégration des mécanismes de notification (e-mail, notifications push via FCM).
            \end{itemize}
        \item \textbf{Durée} : 3 semaines.
        \item \textbf{Tâches clés} :
            \begin{itemize}
                \item Développement du module de gestion des alertes.
                \item Intégration des notifications push et e-mail.
                \item Tests de validation des alertes et des notifications.
            \end{itemize}
        \item \textbf{Dépendances} : Les données collectées dans les Sprints 1 et 2 doivent être disponibles pour ce sprint.
        \item \textbf{Risques} :
            \begin{itemize}
                \item Retards dans la livraison des alertes critiques.
                \item Solution : Optimisation du temps de réponse et tests de performance.
            \end{itemize}
    \end{itemize}
\end{itemize}